%! Modeling with Exponential Functions: Interest, Half-Life, and e — Unit 6, Lesson 6.4 — Lesson Plan
\documentclass[10pt]{article}
\usepackage{algebra2-article}
\usepackage{algebra2-boxes}
\usepackage{graphicx}   % -article does NOT load graphicx; needed for thumbnails/screenshots
\graphicspath{{images/}}
\newcommand{\TallMath}[1]{$\displaystyle #1\rule[-1.4em]{0pt}{3.2em}$}
\newcommand{\CourseName}{Algebra 2}
\newcommand{\MeetingLength}{60 minutes}
\newcommand{\UnitNumberName}{Unit 6: Exponential Functions \quad}
\newcommand{\LessonNumberName}{Lesson 6.4: Modeling --- Compound Interest, Half-Life, and $e$}

\begin{document}
\begin{center}
    {\Large\bfseries \CourseName \\
    {\small\bfseries \UnitNumberName \LessonNumberName}}
\end{center}

\begin{tcolorbox}[colback=forestbg, colframe=forest, boxrule=0.9pt, arc=2mm,
    left=2mm, right=2mm, top=1.5mm, bottom=1.5mm]
    \textbf{Primary Objective:} Students can build, evaluate, and interpret exponential models drawn
    from money and medicine --- compound interest $A=P\left(1+\frac{r}{n}\right)^{nt}$, continuous
    growth $A=Pe^{rt}$, half-life $A_{0}\left(\frac12\right)^{t/h}$, and doubling
    $A_{0}2^{\,t/d}$ --- and can say, for any of them, what one period is and how many periods the
    exponent is counting. Students use the \textbf{range} and \textbf{horizontal asymptote} to state
    what a decay model never does, and they leave able to distinguish an equation with \emph{no}
    solution from one whose solution exists but is out of the reach of a common base.
    \\[2pt]
    \textbf{Standards (2023 VA SOL): A2.F.2f} (primary --- represent and interpret contextual
    situations with exponential functions, including evaluating and explaining the meaning of the
    result), with \textbf{A2.F.2a} (domain and range, including the range boundary a decay model never
    crosses) and \textbf{A2.F.2h} (horizontal asymptotes --- the standard names exponential functions
    explicitly). \textbf{A2.F.1e} is exercised throughout, since each model is written from a verbal
    description and then read back from a table or a graph. \emph{Note:} as documented in 6.3, the
    2023 standards contain no \texttt{A2.EI} strand for exponential equations, so the closing wall is
    taught as \emph{modeling}, not as solving --- the equation is set up and located on a graph, and
    left unsolved on purpose.
\end{tcolorbox}

\renewcommand{\arraystretch}{1.4}
\begin{skillbox}[Priority Ideas \& Skills]{goldbox}
\begin{tabularx}{\linewidth}{@{}X|X@{}}
\textbf{Priority Ideas \& Skills} & \textbf{Key Understandings (the ``why'')} \\[2pt]
\hline
\vspace{-6pt}
\begin{itemize}[leftmargin=1.1em, nosep, after=\vspace{2pt}]
  \item Identify $P$, $r$, $n$, and $t$ from a verbal description and substitute into $A=P\left(1+\frac{r}{n}\right)^{nt}$.
  \item Convert an annual rate into a rate per period ($\frac{r}{n}$) and a time in years into a count of periods ($nt$).
  \item Read the table $\left(1+\frac1m\right)^{m}$ as closing in on $e\approx2.71828$, and use $A=Pe^{rt}$.
  \item Compare compounding schedules and quantify how little the extra chopping is worth.
  \item Build $A_{0}\left(\frac12\right)^{t/h}$ from a half-life and $A_{0}2^{\,t/d}$ from a doubling time.
  \item Evaluate at a time that is \emph{not} a whole number of periods, and check the answer's size against the table.
  \item State the range and horizontal asymptote of a decay model and interpret them in context.
  \item Set up the doubling equation, bracket its answer on a graph, and say why a common base cannot finish it.
\end{itemize}
&
\vspace{-6pt}
\begin{itemize}[leftmargin=1.1em, nosep, after=\vspace{2pt}]
  \item \emph{The base is what happens in one period; the exponent counts the periods.} Every formula today is that sentence.
  \item Nothing here is a new function family --- it is Lesson 6.1's $y=ab^{x}$ with a shorter ``one period.''
  \item The rate is divided and the time is multiplied \textbf{together}. Doing one without the other is the whole error surface.
  \item Compounding more often has a \textbf{ceiling}, and $e$ is that ceiling. This is why continuous models exist, not a way to earn more.
  \item Half-life and interest are the same structure --- the clock changed, not the mathematics.
  \item A fractional exponent means a time between two periods. It is normal, and Lesson 5.1 already licensed it.
  \item ``It never reaches zero'' is a claim about the asymptote (6.0), and it is where the model stops describing reality.
  \item Knowing exactly which questions this unit cannot answer is what makes Unit 7 necessary rather than arbitrary.
\end{itemize}
\end{tabularx}
\end{skillbox}

\begin{skillbox}[{Vocabulary, Concepts, \& Theorems}]{sky}
\renewcommand{\arraystretch}{1.3}
\begin{tabularx}{\linewidth}{@{}l X@{}}
\textbf{Principal, $P$} & The starting amount, before any interest. It is the $a$ of $y=ab^{x}$. \\
\textbf{Compounding period} & The time between two moments when interest is added; $n$ of them fit in a year, so each earns $\frac{r}{n}$. \\
\textbf{Compound interest} & \TallMath{A = P\left(1+\frac{r}{n}\right)^{nt}} \\
\textbf{The number $e$} & $e\approx2.71828$, the value $\left(1+\frac1m\right)^{m}$ closes in on as $m$ grows. Irrational, like $\pi$. \\
\textbf{Continuous compounding} & \TallMath{A = Pe^{\,rt}} --- the ceiling: no $n$ is left to increase. \\
\textbf{Half-life $h$} & The time for half of a quantity to disappear: $A(t)=A_{0}\left(\frac12\right)^{t/h}$. \\
\textbf{Doubling time $d$} & The time for a quantity to double: $A(t)=A_{0}\,2^{\,t/d}$. \\
\textbf{Why $t/h$, not $t$} & The exponent must be counted in \emph{periods}. Dividing by $h$ converts hours (or years) into halvings. \\
\textbf{Out of reach vs.\ nonexistent} & $(1.05)^{t}=2$ has an answer no common base can name; $2^{x}=-3$ has none at all. \\
\end{tabularx}
\end{skillbox}

\begin{fixedskillbox}[Activate Prior Knowledge \& Spiral Review]{forestbg}
    The Warm-Up rehearses exactly what today needs, in the order the lesson uses it. \textbf{(1)} A
    $6\%$ annual rate split across $n=1,4,12$: rate per period, factor per period, and the number of
    periods in five years. This is Lesson 6.1's percent-to-factor translation (\texttt{A2.F.2f}) with
    one new column, and \emph{both} of the day's signature errors are the quarterly row done half-way.
    Debrief it at the board and leave $0.015$, $1.015$, $20$ written where they will stay all period.
    \textbf{(2)} Powers of $\frac12$, including $\left(\frac12\right)^{1/2}$ and
    $\left(\frac12\right)^{1.5}$ (Unit 5 rational exponents), then the question that matters: how many
    \emph{halvings} happen in 18 hours, and in 9 hours? The answer $1.5$ is the half-life exponent
    arriving before the formula does --- do not resolve it algebraically, just point at the decimal
    they already computed. \textbf{(3)} Lesson 6.3, twice: solve $2^{t}=32$ by common base, then look
    at $(1.05)^{t}=2$ and say why Step 1 is impossible. Accept a short answer; Section 4 finishes it.
    If a student says ``no solution,'' correct it immediately --- that distinction is the spine of the
    last ten minutes.
\end{fixedskillbox}

\begin{skillbox}[Hook]{forestbg}
    Four banks, printed on the page, all advertising $6\%$ on the same \$1000 for one year: Ashmore
    compounds annually and pays \$1060.00; Bellhaven quarterly, \$1061.36; Corliss monthly, \$1061.68;
    Dunmore daily, \$1061.83. \textbf{Do not compute these at the board} --- the surprise is not the
    arithmetic, it is the differences. Ask: \textbf{``Dunmore stops and adds interest 365 times as often
    as Ashmore. What does all that extra work earn you?''} The answer, \$1.83 over an entire year, lands
    hard, because nearly every student predicts that daily compounding is dramatically better. Then make
    them look at the pattern of gains: $+\$1.36$, then $+\$0.32$, then $+\$0.15$. \textbf{``Each extra
    chop of the year buys less than the one before. Where does that stop?''} Close by taking a real
    prediction and writing the guesses on the board: \textbf{is there a bank so aggressive --- compounding
    every second, every instant --- that it doubles your money in a year, or is something stopping all of
    them?} Section 2 settles it, and the answer has a name.
\end{skillbox}

\begin{skillbox}[Lesson]{forestbg}
\begin{multicols}{2}
\textbf{1. Cutting the year into pieces.} Start from Lesson 6.1's $y=ab^{x}$ and change exactly one
thing: ``one period'' no longer has to be a year. Build $\frac{r}{n}$ and $nt$ as a pair, never
separately, then box $A=P\left(1+\frac{r}{n}\right)^{nt}$ and read it aloud as the sentence the whole
lesson runs on: \emph{the base is what happens in one period, and the exponent counts the periods.}
Work \$2500 at $4\%$ quarterly for 6 years line by line --- $P$, $r$, $n$, $t$, then $\frac{r}{n}=0.01$,
then $nt=24$, then $A=\$3174.34$ --- and set it beside the annual answer, \$3163.30. The extra \$11.04
is deliberately unimpressive; it is the Hook restated in a second context.

\textbf{2. Where the gains stop.} Strip the money out and ask the cleanest version: \$1 at $100\%$ for
a year, in $m$ pieces, is $\left(1+\frac1m\right)^{m}$. Fill the last three cells together and say what
students can see --- the digits are freezing from the left. Name $e\approx2.71828$ as \emph{the number
the chopping runs into}, not as a formula, and put it next to $\pi$: irrational, endless, and on the
calculator. Then $A=Pe^{rt}$, with the observation that the $n$ is gone because there is nothing left to
increase. Return to the worked example: continuously, \$2500 becomes \$3178.12 --- \$3.78 more than
quarterly. That is the answer to the Hook, and it should be said as a sentence: \textbf{compounding
schedules are not where returns come from.}

\columnbreak

\textbf{3. A different clock.} Half-life and doubling time are the same structure, so present them as
a pair and let the sentence do the work again: the base says what one period does, and $t/h$ counts the
periods. \textbf{Spend real time on why the exponent is $t/h$ and not $t$} --- it is a unit conversion,
hours into halvings, and the Warm-Up already produced $1.5$. Fill the 80 mg tracer table across, saving
$t=9$ for last: a fractional exponent, an answer of $28.28$ mg, and a size check against $40$ and $20$.
Then the Lesson 6.0 question --- will it ever be gone? --- answered from the \textbf{asymptote} and the
\textbf{range} $(0,80]$, never from arithmetic. Close the section with the green box: $10$ mg \emph{is}
reachable by yesterday's method, because $10$ is $80$ halved three times.

\textbf{4. The wall.} Ask the question every investor asks: when does \$1000 at $5\%$ double? Set up
$1000(1.05)^{t}=2000$, divide to $(1.05)^{t}=2$, and go to the graph. The gold line crosses, between
$t=14$ (\$1979.93) and $t=15$ (\$2078.93), so \textbf{the answer exists}. Then run Step 1 of Lesson 6.3
and watch it fail: $2$ is not a power of $1.05$. Put that beside $2^{x}=-3$, which has no solution at
all, and insist on the difference in words. Post the running list --- now five equations long. Do not
name logarithms.
\end{multicols}
\end{skillbox}

\begin{skillbox}[Active Monitoring]{redbox}
Circulate during the notes practice and Tier work. Watch for and correct:
\begin{itemize}[leftmargin=1.2em, nosep]
  \item the annual rate kept as the base while the exponent is converted --- $1000(1.08)^{20}$ for a quarterly account; the tell is an absurd balance;
  \item the rate split but the time not --- $1000(1.02)^{5}$; the tell is a balance barely above the principal;
  \item $t$ used where $t/h$ belongs in a half-life model --- $90\left(\frac12\right)^{18}$ instead of $90\left(\frac12\right)^{18/6}$;
  \item $r$ entered as a percent rather than a decimal ($n=4$, $r=8$), which is Lesson 6.1's error one unit later;
  \item rounding the base or an intermediate power before the final step --- money answers drift by dollars;
  \item distrusting a fractional number of half-lives and ``fixing'' $t/h=3.4$ to $3$;
  \item answering ``the drug is gone'' or ``the truck is worthless'' without the asymptote --- the Unit 5 endpoint habit that 6.0 was built to break;
  \item \textbf{answering ``no solution'' for $(1.05)^{t}=2$}; the correct statement is that the solution exists and today's method cannot reach it;
  \item in Tier E, treating the Rule of 72 as exact rather than as an estimate calibrated near $8\%$.
\end{itemize}
Cold-call prompts: ``What is one period here?'' ``How many of them fit in the time you were given?''
``Did you divide the rate \emph{and} multiply the time, or only one?'' ``Is that answer a reasonable
size --- would \$5000 really become \$82{,}000 in four years?'' ``Does that equation have no answer, or
an answer we can't reach yet? Those are different sentences.''
\end{skillbox}

\begin{skillbox}[Group Work \& Differentiation]{redbox}
\textbf{The Base Is One Period. The Exponent Counts Periods.} activity (tiered; mirrors the notes).
\begin{multicols}{3}
\textbf{Tier R --- Remediate.} A setup-before-you-solve table for \$2000 at $3\%$ over 10 years,
annually versus quarterly (the gain is \$8.87 --- ask whether that matched the group's guess); a 600 mg
medicine with a 12-hour half-life, filled across four halvings, closing on \emph{which feature of the
graph says it never reaches zero}; and a three-row diagnostic in which students name \emph{what one
period is} for a quarterly model, a half-life model, and a doubling model. \textbf{Check Part 3 first
when you circulate} --- it predicts everything else on the page.

\columnbreak
\textbf{Tier A --- Approaching.} \$5000 at $3.6\%$ for eight years, monthly (\$6665.91) against
continuously (\$6668.79), with a written response to the classmate who expects continuous compounding to
win by a lot; a 120 mg dose with a 5-hour half-life evaluated at $t=10$ and at $t=17$, where $3.4$
halvings must be justified and size-checked against $15$ and $7.5$ mg; and the \textbf{error analysis}
--- Renata writes $1000(1.08)^{20}$, Miles writes $1000(1.02)^{5}$. They are opposite halves of one
mistake and need different re-teaches, which is exactly what the closing question asks students to
articulate.

\columnbreak
\textbf{Tier E --- Extension.} Three parts. \emph{The ceiling, at a real rate:}
$\left(1+\frac{0.06}{n}\right)^{n}$ across $n=1,4,12,365$ and then $e^{0.06}=1.061837$, with ``how much
more could compounding every \emph{second} pay?'' (about half a cent). \emph{The Rule of 72:} tested at
$5\%$, $8\%$, and $12\%$ against the true doubling times --- it is excellent near $8\%$ and drifts on
either side, which is calibration, not a formula. Students who ask \emph{why 72} are asking a logarithm
question; say so and stop. \emph{Which can you finish?} Four modeling questions reduced to equations,
two solvable by common base and two not --- including the same drug and half-life targeted at $75$ mg
(works) and $200$ mg (does not).
\end{multicols}
\end{skillbox}

\begin{skillbox}[Individual Work \& Assessment]{redbox}
\textbf{Exit Ticket} (independent, no notes), four items on one page: \textbf{(1)} \$4000 at $3\%$
compounded monthly for 5 years, with the base ($1.0025$) and exponent ($60$) written before the answer
(\$4646.47); \textbf{(2)} a 50 mg dose with a 4-day half-life --- write the model, evaluate at 12 days
($6.25$ mg), and say from the \emph{asymptote} whether it ever reaches $0$; \textbf{(3)} an
\textbf{SOL-style multiple-choice item} --- \$1500 at $4.8\%$ compounded quarterly for 7 years, with
distractors $1500(1.048)^{7}$ (neither converted), $1500(1.012)^{7}$ (Miles's error), and
$1500(1.048)^{28}$ (Renata's error) against the correct $1500(1.012)^{28}$, so the wrong answer names
the re-teach; and \textbf{(4)} the wall: $(1.012)^{4t}=2$ cannot be finished with a common base ---
\emph{does it still have an answer?} ``No'' is wrong, and it is the single most informative response on
the page. Sort the collected tickets into \emph{got it / setup / asymptote / the wall}. The
\emph{setup} pile needs practice, not re-teaching. A large \emph{wall} pile changes how 6.5 opens and
must be fixed before Unit 7: five minutes on the Section 4 graph with the gold line visibly crossing.
\end{skillbox}

\begin{skillbox}[Reinforcement \& Extension]{goldbox}
\textbf{Homework:} four accounts at four compounding frequencies, with the base and exponent written
before the balance (the daily row's exponent of $1460$ is the one to check); \$6000 at $4\%$ for ten
years quarterly against continuously, whose \$17.77 gap must be \emph{explained}; a caffeine model
(200 mg, 5-hour half-life) evaluated at a non-half-life time and size-checked against the table, then
asked whether the caffeine is ever gone; a \textbf{graph read} of $800(1.06)^{t}$ against $y=1600$,
where students bracket the doubling between years 11 and 12, write the equation, and say what stops
Step 1; and \textbf{error analysis} repeating the day's two errors in two different models --- Tobias
keeps an annual base with a monthly count, Marisol counts hours instead of halvings, and both absurd
answers (\$81{,}969.36 from \$5000; $0.00034$ mg from 90 mg) are reasonableness checks students should
have run themselves. \textbf{Extension:} a \$32{,}000 truck losing $18\%$ a year, where depreciation and
half-life are revealed to be the same object --- the base is $0.82$ and not $0.18$, the ``half-life of a
truck'' is $(0.82)^{t}=0.5$ (between years 3 and 4, no common base, fifth entry on the running list),
and the model's asymptote says the truck is never worthless, which is mathematically right and
physically wrong. That last beat is deliberate set-up for 6.5's extrapolation discussion.
\textbf{Preview:} Lesson~6.5, \emph{Exponential Regression \& Choosing a Model} --- where nobody hands
you the rate, and Lesson 6.0's fingerprint test decides the family before technology fits anything.
\end{skillbox}
\end{document}
